% Einheit 3 von 5 – Brüche vergleichen (bank/brueche-dezimalzahlen/e3.jsonl, zusammenbau v0.1)
%% TODO Zweigzeile Teil 1: was hier gelernt wird (Satz in Schülersprache aus dem Einheitstitel) – Einheit 3
\einheitenkopf[e3]{Einheit 3 von 5 · Brüche vergleichen}
\zweigzeile{ab Kl. 5, je nach Buch bis Kl. 6 · P10 oft}
\setcounter{aufgabe}{17}

%% TODO Titel: Ich-kann-Satz fehlt in der Bank (Platzhalter: Kettenname) – Nr. 18
%% TODO Anweisung: ein Satz über der ersten Teilaufgabe fehlt in der Bank – Nr. 18
\begin{aufgabe}{Vergleichen}
\begin{teile}
\teil $\frac{3}{8}$ und $\frac{5}{8}$ – welcher Vergleichsweg ist am schnellsten? Kreuze an.\\ \kreuz{gleicher Nenner}\\ \kreuz{gleicher Zähler}\\ \kreuz{mit einem Halben vergleichen}\\ \kreuz{gleichnamig machen}
\end{teile}
%% TODO Darstellung für schwach fehlt in der Bank (brueche-dezimalzahlen-e3-k1-s1-v1; Abschnitt „Für schwache Schüler“)
\swz{Welcher Bruch ist größer: $\frac{2}{11}$ oder $\frac{7}{11}$?}{}
%% TODO Darstellung für schwach fehlt in der Bank (brueche-dezimalzahlen-e3-k1-s1-v2; Abschnitt „Für schwache Schüler“)
\swz{Welcher Bruch ist größer: $\frac{7}{12}$ oder $\frac{5}{12}$?}{}
%% TODO Darstellung für schwach fehlt in der Bank (brueche-dezimalzahlen-e3-k1-s1-v3; Abschnitt „Für schwache Schüler“)
\swz{Welcher Bruch ist größer: $\frac{1}{6}$ oder $\frac{5}{6}$?}{}
%% TODO Darstellung für schwach fehlt in der Bank (brueche-dezimalzahlen-e3-k1-s1-v4; Abschnitt „Für schwache Schüler“)
\swz{Welcher Bruch ist größer: $\frac{11}{15}$ oder $\frac{8}{15}$?}{}
%% TODO Darstellung für schwach fehlt in der Bank (brueche-dezimalzahlen-e3-k1-s1-v5; Abschnitt „Für schwache Schüler“)
\swz{Welcher Bruch ist größer: $\frac{9}{20}$ oder $\frac{13}{20}$?}{}
\swfrage{Was bleibt gleich, was ändert sich?}
%% TODO Darstellung für schwach fehlt in der Bank (brueche-dezimalzahlen-e3-k1-s2-v1; Abschnitt „Für schwache Schüler“)
\swz{Welcher Bruch ist größer: $\frac{3}{8}$ oder $\frac{3}{5}$?}{}
%% TODO Darstellung für schwach fehlt in der Bank (brueche-dezimalzahlen-e3-k1-s3-v1; Abschnitt „Für schwache Schüler“)
\swz{Welcher Bruch ist größer: $\frac{4}{9}$ oder $\frac{5}{8}$? Vergleiche mit $\frac{1}{2}$.}{}
%% TODO Darstellung für schwach fehlt in der Bank (brueche-dezimalzahlen-e3-k1-s4-v1; Abschnitt „Für schwache Schüler“)
\swz{Welcher Bruch ist größer: $\frac{3}{4}$ oder $\frac{5}{7}$?}{}
\swa{Trage $\frac{3}{8}$ und $\frac{7}{8}$ am Zahlenstrahl ein.}{\zahlenstrahl[xmin=0,xmax=1,xstep=0.125,karo=1.2]{0/0, 1/1}}
%% TODO Darstellung für schwach fehlt in der Bank (brueche-dezimalzahlen-e3-k1-s6-v1; Abschnitt „Für schwache Schüler“)
\swz{Ordne von klein nach groß: $\frac{2}{3}$, $\frac{1}{6}$, $\frac{5}{6}$.}{}
\swz[3]{Gib eine Zahl an, die zwischen $\frac{1}{4}$ und $\frac{3}{5}$ liegt. \hfill (P10 2014 OS)}{}
\end{aufgabe}

%% TODO Titel: Ich-kann-Satz fehlt in der Bank (Platzhalter: Kettenname) – Nr. 19
%% TODO Anweisung: ein Satz über der ersten Teilaufgabe fehlt in der Bank – Nr. 19
\begin{aufgabe}{Zahl zwischen zwei Brüchen angeben}
%% TODO Darstellung für schwach fehlt in der Bank (brueche-dezimalzahlen-e3-k2-s1-v1; Abschnitt „Für schwache Schüler“)
\swz{Gib einen Bruch zwischen $\frac{1}{5}$ und $\frac{2}{5}$ an.}{}
\end{aufgabe}

%% TODO Titel: Ich-kann-Satz fehlt in der Bank (Platzhalter: Kettenname) – Nr. 20
%% TODO Anweisung: ein Satz über der ersten Teilaufgabe fehlt in der Bank – Nr. 20
\begin{aufgabe}{Vergleichen – Fehler finden, Begründen, Darstellungswechsel, Anwendung}
\swz[3]{Kai sagt: $\frac{2}{3}$ und $\frac{5}{6}$ sind gleich groß, denn bei beiden fehlt nur ein Stück zum Ganzen. Finde den Fehler.}{}
\swfrage{Warum ist $\frac{1}{5}$ größer als $\frac{1}{6}$? Begründe ohne Rechnung.}
\swz{Welcher Bruch gehört zu A?}{\zahlenstrahl[xmin=0,xmax=1,xstep=0.25,karo=3]{0/0, 1/1, 0.75/A}}
\swz[3]{Lara und Emil laden gleich große Dateien. Bei Lara sind $\frac{5}{8}$ geladen, bei Emil $\frac{2}{3}$. Wer ist schon weiter?}{}
\end{aufgabe}

\uebersichtskasten{Gleicher Nenner: größerer Zähler, größerer Bruch. \quad 5/9 > 4/9 \\ Gleicher Zähler: größerer Nenner, kleinerer Bruch (feinere Teile). \quad 2/5 > 2/7 \\ Sonst: mit 1/2 oder 1 vergleichen – oder gleichnamig machen. \quad 3/7 < 1/2 < 5/9 \quad 2/3 und 3/5: 10/15 > 9/15 \\ Zwischen zwei Brüchen liegt immer noch ein Bruch: erweitern, dann dazwischen wählen. \quad zwischen 3/10 und 4/10 liegt 35/100}
